% region
    \ptitle{Ecuații cu numere reale. Expresii algebrice elementare}
        \textbf{\underbar{Def.}} O expresie algebrică elementară este însușită de numere reale și litere împreună de 
        operațile elementare, ele fiind adunarea, înmulțirea, împărțirea, ridicarea la putere și radicalii. \\
        O expresie algebrică se notează cu în următorul mod.
        \[ E(x, y, z, a, b, c, ...) = ... \text{ sau } F(x, y, z, ...) = ... \text{ sau orice literă} \]
        \textbf{Exemple:}
        \begin{enumerate}
            \item $E(x) = ax + b, \quad  a, b \in \mathbb{R} \text{ și } a \neq 0$ - forma generală a expresiei de gradul întai \\
                \textbf{\underbar{Obs.}} Daca egalăm expresia de gradul întai cu 0, obținem ecuația de gradul întai.
                \[ E(x) = 0 \iff ax + b = 0, \quad a, b \in \mathbb{R} \text{ și } a \neq 0 \text{ - ecuația de gradul întai} \]
            \item $E(x) = ax^2 + bx + c, \quad  a, b, c \in \mathbb{R} \text{ și } a \neq 0$ - forma generală a expresiei de gradul al doilea
                \[ E(x) = 0 \iff ax^2 + bx + c = 0, \quad a, b, c \in \mathbb{R} \text{ și } a \neq 0 \text{ - forma generală a expresiei de gradul al doilea}\]
            \item $E(x) = a_0 x^2 + a_1 x^2y + a_2 x^2 y^2 + a_3 x y^2 + a_4 y^2 + b_0 x + b_1 xy + b_2 y + c,\quad a_n, b_m, c \in \mathbb{R}, \forall n, m \in \mathbb{N} \text{ și } \exists n \in \mathbb{N} \text{ a.î. } a_n \neq 0$ - forma generală a expresiei
                de gradul al doilea cu două necunoscute
        \end{enumerate}
        Valoarea unei expresii la o altă valoare sau o pereche de valori, noi înlocuim necunoscuta cu valoarea sau perechea de valori în ordine. \\
        \textbf{Exemplu:}
        \[ E(x, y) = x^2 + y\sqrt{x} + 2y,\ \forall x \in \mathbb{R}_+ \cup \{0\} = [0, \infty), y \in \mathbb{R} \implies E(4, -1) = 4^2 - \sqrt{4} - 2 = 16 - 2 - 2 = 12 \]
        Partea unei expresii care este separată de semnele plus sau minus de altele se numește un termen, ea fiind formată dintr-un coeficient constant și o parte 
        dependentă de necunoscute.
        \[ E(x) = ax^2 + bx + c \implies (ax^2, bx, c) \text{ - termeni} \text{ și } (a, b, c) \text{ - coeficienți}\]
        \psubtitle{Adunarea expresiilor algebrice elementare}
        Adunăm două expresii algebrice elementare utilizând principiul factorului comun ($ax + bx = (a + b)x, \ \forall a, b, x \in \mathbb{R}$) în gruparea termenilor
        asemenea, adunând coeficienții. \\
        \textbf{Exemplu:}
        \[ (5x^2 - 3x + 6) + (-x^2 + x - 2) = (5 - 1)x^2 + (-3 + 1)x + 6 - 2 = 4x^2 - 2x + 4 \]